%% 2i2d-poutre-deux-appuis — poutre de 6,00 m articulée en A, sur appui glissant
%% en B, portant deux charges ponctuelles. Les réactions R_A et R_B sont les inconnues.
\documentclass[tikz,border=4pt]{standalone}
\usepackage{liaisons}
\begin{document}
\begin{tikzpicture}[x=1cm,y=1cm,
  force/.style={-{Stealth[length=5pt]}, line width=1pt, solideA},
  inconnue/.style={-{Stealth[length=5pt]}, line width=1pt, solideF},
  cote/.style={{Stealth[length=4pt]}-{Stealth[length=4pt]}, line width=0.6pt, solideD},
  attache/.style={line width=0.4pt, solideD!80},
  axe/.style={-{Stealth[length=5pt]}, line width=0.6pt, solideF},
  etiq/.style={font=\small, inner sep=2pt},
  petit/.style={font=\scriptsize, inner sep=1.5pt}]

\newcommand\batihoriz[3]{%
  \draw[line width=1pt, solideE] (#1,#3) -- (#2,#3);
  \begin{scope}
    \clip (#1,#3-0.24) rectangle (#2,#3);
    \foreach \hx in {-16,-15.76,...,16} { \draw[line width=0.5pt, solideE!75] (\hx,#3) -- (\hx-0.28,#3-0.28); }
  \end{scope}}
\newcommand\appuitri[3]{%
  \draw[line width=1pt, draw=solideE, fill=solideE!12, line join=round]
    (#1,#2) -- (#1-0.3,#3) -- (#1+0.3,#3) -- cycle;}

%% --- la poutre --------------------------------------------------------------
\draw[line width=2.6pt, solideB, line cap=round] (0.9,1.6) -- (6.9,1.6);

%% --- appui A : articulation ---------------------------------------------------
\appuitri{0.9}{1.57}{1.15}
\batihoriz{0.45}{1.35}{1.15}
\node[etiq, anchor=north] at (0.9,0.88) {$A$};

%% --- appui B : appui glissant (triangle + rouleaux) -----------------------------
\appuitri{6.9}{1.57}{1.22}
\draw[line width=0.8pt, draw=solideE, fill=white] (6.72,1.11) circle (0.11);
\draw[line width=0.8pt, draw=solideE, fill=white] (7.08,1.11) circle (0.11);
\batihoriz{6.45}{7.35}{1.0}
\node[etiq, anchor=north] at (6.9,0.88) {$B$};

%% --- les charges ----------------------------------------------------------------
\draw[force] (2.9,2.6) -- (2.9,1.68);
\node[etiq, anchor=south, text=solideA] at (2.9,2.62) {12 kN};
\draw[force] (4.9,2.6) -- (4.9,1.68);
\node[etiq, anchor=south, text=solideA] at (4.9,2.62) {6 kN};

%% --- les réactions d'appui (inconnues) ---------------------------------------------
\draw[inconnue] (0.9,1.66) -- (0.9,2.5);
\node[etiq, anchor=east, text=solideF] at (0.82,2.2) {$R_A$};
\draw[inconnue] (6.9,1.66) -- (6.9,2.5);
\node[etiq, anchor=west, text=solideF] at (6.98,2.2) {$R_B$};

%% --- les cotes -------------------------------------------------------------------------
\foreach \x in {0.9,2.9,4.9,6.9} { \draw[attache] (\x,0.8) -- (\x,0.0); }
\draw[cote] (0.9,0.45) -- (2.9,0.45);
\draw[cote] (2.9,0.45) -- (4.9,0.45);
\draw[cote] (4.9,0.45) -- (6.9,0.45);
\foreach \x in {1.9,3.9,5.9} { \node[petit, anchor=south, text=solideD] at (\x,0.48) {2,00 m}; }
\draw[cote] (0.9,0.05) -- (6.9,0.05);
\node[petit, anchor=south, text=solideD] at (3.9,0.08) {6,00 m};

%% --- repère ------------------------------------------------------------------------------
\draw[axe] (0.15,0.05) -- ++(0.55,0) node[petit, anchor=north] {$\vec{x}$};
\draw[axe] (0.15,0.05) -- ++(0,0.55) node[petit, anchor=east] {$\vec{y}$};
\end{tikzpicture}
\end{document}
